\begin{textsample}{NEG Dim 4 – 2023\_c6 – Score -8.00 – 2023\_c6/enade\_2023\_med\_c56e7c1ea6.txt}  \label{ex:f4_neg_018}
CAPÍTULO 5 ANÁLISE TÉCNICA DA PROVA Este capítulo tem o objetivo de apresentar o desempenho dos estudantes concluintes de Medicina no Enade 2023.
Para isso, foram calculadas as estatísticas básicas da prova em seu todo ( subseção 5.1.1 ), bem como as estatísticas relacionadas aos componentes de Formação Geral ( subseção 5.1.2 ) e de Conhecimento Específico da área ( subseção 5.1.3 ).
Dadas as suas características, foram analisadas, em separado, as questões objetivas ( seção 5.2 ) e as questões discursivas ( seção 5.3 ).
Tomando -se como base a questão discursiva do componente de Formação Geral, nas subseções de 5.3.1.4 a 5.3.1.6, são apresentados os resultados e os comentários sobre a correção das respostas em relação ao desempenho linguístico.
Para as questões objetivas, foram disponibilizados os índices de facilidade e de discriminação ponto-bisserial, também separadamente, para os componentes de Formação Geral ( subseção 5.2.1 ) e de Conhecimento Específico ( subseção 5.2.2 ).
De cada componente, uma das questões foi escolhida para exemplificar a análise gráfica, relacionando -se as alternativas escolhidas pelos estudantes ( inclusive o gabarito ) ao número de acertos no componente.
No Anexo IV, apresenta -se a íntegra da análise gráfica para todas as questões objetivas.
Para cada uma das questões discursivas, os conteúdos dos tipos mais comuns de respostas dos estudantes são apresentados e comparados com o padrão de resposta esperado ( ver Anexo VIII ).
Nas tabelas, constam as seguintes estatísticas das notas5 : média do desempenho na prova, erro-padrão da média, desvio-padrão, nota mínima, mediana e nota máxima para cada um de seus componentes.
Tais estatísticas contemplam o total de estudantes concluintes da área de Medicina inscritos e presentes na prova do Enade 2023, tendo em vista agregações, ou por grandes regiões e o país como um todo, ou por categoria administrativa, organização acadêmica da IES e modalidade de oferta dos cursos.
Em relação aos gráficos de distribuição de notas, o intervalo considerado foi de 10 unidades, aberto à esquerda e fechado à direita, com exceção do primeiro intervalo, [ 0 ; 10 ], fechado em ambos os extremos.
Para os gráficos de distribuição das notas das questões discursivas, foram consideradas mais duas categorias : questão em branco6 e nota zero. 5 Uma definição dessas estatísticas pode ser encontrada no Glossário inserido no final desse relatório. 6 Nesse grupo estão incluídas também as respostas classificadas como nulas ou desconsideradas. 5.1 ESTATÍSTICAS BÁSICAS DA PROVA Esta seção apresenta estatísticas selecionadas e histogramas da nota geral ( 5.1.1 ) e de cada componente : Formação Geral ( 5.1.2 ) e Conhecimento Específico ( 5.1.3 ).
São também apresentadas estatísticas selecionadas de subpopulações, caracterizadas por grande região, categoria administrativa, organização acadêmica e modalidade de oferta. 5.1.1 ESTATÍSTICAS BÁSICAS GERAIS Na Tabela 5.1, são apresentadas as estatísticas básicas da prova, por grande região, dos estudantes concluintes de Medicina.
Como se pode verificar, a partir dos dados mostrados na Tabela 5.1, a média das notas da prova foi de 65,0 ( nas seções seguintes serão analisados os componentes de Formação Geral e de Conhecimento Específico ).
Os estudantes da região Norte obtiveram a média mais baixa ( 62,1 ), e os da região Sul alcançaram a média mais alta ( 67,5 ).
As demais médias foram 65,3, na região Centro-Oeste ; 64,5, na região Nordeste ; 64,8, na região Sudeste.
O desvio-padrão para o Brasil foi de 12,2, sendo o maior encontrado na região Sudeste ( 12,4 ), e o menor, na região Norte ( 11,6 ).
A região que obteve a maior nota máxima foi a região Sudeste ( 99,0 ), ao passo que a região que atingiu a menor nota máxima foi o Norte ( 93,3 ).
A mediana do Brasil, como um todo, foi igual a 65,4, sendo a maior mediana obtida na região Sul ( 68,0 ), e a menor, na região Norte ( 62,6 ).
Considerando -se as notas segundo grande região, observa -se que existe diferença estatisticamente significativa ao nível de 95 \% entre as médias obtidas em todas as regiões, exceto entre Sudeste e Centro-Oeste ; Sudeste e Nordeste.
O comportamento das notas dos estudantes de todo o Brasil pode ser observado no Gráfico 5.1.
O Gráfico 5.1 apresenta um histograma com a distribuição das notas dos estudantes de Medicina de todo o Brasil : uma distribuição unimodal com a moda no intervalo ( 60 ; 70 ].
Na Tabela 5.2, são apresentadas informações referentes à média da nota final, desagregadas por categoria administrativa, organização acadêmica e modalidade de oferta dos cursos.
Os dados da Tabela 5.2 mostram que, em relação à categoria administrativa, os estudantes das IES públicas obtiveram média mais alta ( 68,7 ) que a média das IES privadas ( 63,7 ).
O desvio-padrão para as IES públicas foi de 11,9, indicando uma dispersão baixa das notas nesta categoria administrativa.
Já o desvio-padrão da IES privadas foi 12,1.
No que se refere à organização acadêmica, os estudantes de universidades obtiveram a média de 65,6, enquanto os de centros universitários e faculdades obtiveram, respectivamente, médias de 63,6 e 64,9.
Constata -se que existe diferença estatisticamente significativa, ao nível de 95 \%, entre as médias das organizações acadêmicas. 5.1.2 ESTATÍSTICAS BÁSICAS NO COMPONENTE DE FORMAÇÃO GERAL Na Tabela 5.3, são apresentadas as estatísticas básicas em relação ao componente da prova que avalia a Formação Geral dos estudantes concluintes.
A nota de Formação Geral é considerando -se os acertos nas nove questões objetivas e a nota média obtida na questão discursiva, sendo que a nota dessa questão obtida leva em conta 80 \% da nota relativa ao conteúdo e 20 \% da nota relativa ao desempenho linguístico. nas regiões, exceto entre as regiões Nordeste e Centro-Oeste ; Sudeste e Centro-Oeste ; Sudeste e Nordeste.
No Gráfico 5.2, é possível verificar a avaliação do desempenho dos estudantes no componente de Formação Geral, a partir do histograma da distribuição das notas correspondentes.
Nesse gráfico, a distribuição é unimodal, com moda em ( 60 ; 70 ], mas com pouca diferença para o intervalo que o precede ( 50:60 ].
Na Tabela 5.4, são apresentadas as informações referentes ao desempenho dos concluintes no componente de Formação Geral, em diferentes agregações : categoria administrativa, organização acadêmica e modalidade de oferta.
A partir dos dados da Tabela 5.4, observa -se que há diferença estatisticamente significativa entre as médias dos tipos de categoria administrativa.
A maior média foi obtida por estudantes de IES pública ( 70,6 ), e a menor, pelos estudantes de IES privada ( 64,1 ).
Em relação às médias por tipos de organização acadêmica, encontra -se : 66,6 para as universidades ; 63,8 para os centros universitários, e 65,6 para as faculdades.
A nota máxima obtida pelas IES privadas ( 100,0 ) foi igual em comparação às IES públicas.
O desvio-padrão menor foi percebido nas IES públicas ; e, entre as organizações acadêmicas, esse valor foi menor para os centros universitários.
Já no que se refere à modalidade de oferta, os cursos presenciais apresentaram média de 65,8. 5.1.3 ESTATÍSTICAS BÁSICAS DO COMPONENTE DE CONHECIMENTO ESPECÍFICO Na Tabela 5.5, são apresentadas as estatísticas básicas referentes ao componente de Conhecimento Específico da área de Medicina.
A nota do componente de Conhecimento Específico leva em conta as notas da parte objetiva e a média das notas obtidas na questão da parte discursiva da área.
De acordo com a Tabela 5.5, a média do desempenho dos estudantes do Brasil, para o componente de Conhecimento Específico da prova, foi de 64,7.
A maior média foi obtida na região Sul ( 67,3 ), e a menor, na região Norte ( 61,7 ).
As demais médias foram : 65,1, na região Centro-Oeste ; 64,3, na região Nordeste ; 64,5, na região Sudeste.
Quanto à variabilidade das notas, o desvio-padrão do Brasil foi 13,6 ; sendo o maior desvio-padrão observado na região Sudeste ( 13,8 ), e o menor, na região Norte ( 12,8 ).
Os demais desvios foram : 13,0, na região Centro-Oeste ; 13,5, na região Nordeste ; 13,3, na região Sul.
A mediana das notas dos estudantes de todo o Brasil foi de 65,0.
A maior mediana ocorreu na região Sul ( 67,9 ), e a menor, na região Norte ( 62,0 ).
As demais medianas foram : 65,6, na região Centro-Oeste ; 64,6, na região Nordeste ; 64,7, na região Sudeste.
A nota máxima do Brasil, como um todo, foi 100,0 ; sendo obtida por, pelo menos, um estudante na região Centro-Oeste e Sul.
As demais notas máximas foram : 96,7, na região Nordeste ; 95,7, na região Norte ; 99,5, na região Sudeste.
A nota mínima nessa parte foi zero em quase todas as regiões, com exceção da região Nordeste com 3,3 \%.
Observa -se que existe diferença estatisticamente significativa entre todas as regiões na análise das médias das notas do componente de Conhecimento Específico, exceto entre as regiões Sudeste e Centro-Oeste ; Sudeste e Nordeste.
Assim como nos Gráficos 5.1 e 5.2, mostrados anteriormente, no Gráfico 5.3, apresentado acima, foi possível ser realizada uma avaliação do desempenho de concluintes de Medicina, em relação ao componente de Conhecimento Específico, com um histograma da distribuição das notas correspondentes.
Esta também é uma distribuição unimodal, e o grupo modal é o ( 60 ; 70 ].
Na Tabela 5.6, é apresentada uma comparação dos resultados em relação à categoria administrativa, à organização acadêmica e à modalidade de oferta, levando -se em conta o desempenho de estudantes do componente de Conhecimento Específico da prova.
De acordo com os dados expostos na Tabela 5.6, no que se refere à organização acadêmica, a maior média foi das universidades ( 65,3 ), seguida das faculdades ( 64,6 ), e dos centros universitários ( 63,5 ).
O maior desvio-padrão foi 13,7 das universidades e das faculdades.
As universidades obtiveram a maior mediana ( 65,7 ).
Já as notas máximas foram : 100,0 nas universidades, 98,9 nos centros universitários e 99,5 nas faculdades.
A nota mínima foi zero para todas as organizações acadêmicas.
Observa -se que existe diferença estatisticamente significativa, ao nível de 95 \%, no componente de Conhecimento Específico entre as notas dos tipos de organização acadêmica.
Quanto à categoria administrativa, existe diferença estatisticamente significativa entre as médias das IES públicas ( 68,1 ) e as das IES privadas ( 63,6 ).
Neste caso, assim como no componente de Formação Geral, a maior média foi obtida por estudantes de IES públicas de ensino.
Os cursos de educação presencial apresentaram uma média de 64,7. 5.2 ANÁLISE DAS QUESTÕES OBJETIVAS Esta seção apresenta estatísticas selecionadas e histogramas das questões objetivas dos componentes de Formação Geral ( 5.2.1 ) e de Conhecimento Específico ( 5.2.2 ).
São, também, apresentadas e comparadas as médias das subpopulações caracterizadas por grande região. 5.2.1 COMPONENTE DE FORMAÇÃO GERAL Na Tabela 5.7, são apresentadas as estatísticas básicas relativas às nove questões objetivas do componente da prova que abrange a Formação Geral dos estudantes.
De acordo com os dados da Tabela 5.7, em relação ao componente de Formação Geral, a média do Brasil foi 69,8.
A menor média foi encontrada na região Norte ( 66,6 ), e a maior, na região Sul ( 72,5 ).
As demais médias foram : 70,0, na região Centro-Oeste ; 69,6, na região Nordeste ; 69,4, na região Sudeste.
O desvio-padrão do Brasil foi 17,9 ; sendo o maior desvio-padrão encontrado na região Nordeste ( 18,2 ), e o menor desvio-padrão encontrado na região Sul ( 17,1 ).
A mediana do Brasil, como um todo, foi 66,7, bem como das regiões Centro-Oeste, Nordeste, Norte e Sudeste, exceto a região Sul que atingiu a maior mediana ( 77,8 ).
Quanto à nota máxima foi 100,0 para todas as regiões.
Já as notas mínimas foi : 0,0 também em todas as regiões.
Na Tabela 5.8, exposta mais adiante, são apresentados o índice de facilidade e o índice de discriminação ( ponto-bisserial ) para cada uma das questões objetivas do componente de Formação Geral.
De acordo com os dados da Tabela 5.8, as questões objetivas do componente de Formação Geral, segundo o índice de facilidade, foram assim avaliadas : duas questões foram classificadas como “ Muito fácil ” ; cinco questões foram classificadas como “ Fácil ” ; uma questão foi classificada como"Médio ” ; e uma questão foi classificada na categoria “ Difícil ”.
O índice de facilidade variou de 0,36 a 0,92.
Já em relação ao índice de discriminação, sete questões foram classificadas com o índice “ Muito bom ” ; duas com o índice “ Bom ”, nenhuma questão com o índice “ Médio ” ; e nenhuma com o índice"Fraco ”.
O índice de discriminação variou de 0,37 a 0,47.
No Gráfico 5.4, para exemplificar, analisa -se o comportamento da Questão Objetiva 5 de Formação Geral.
Trata -se de uma questão considerada com nível “ Fácil ” de dificuldade, em relação à facilidade, e que obteve o índice de discriminação 0,46 ( muito bom ).
Na Tabela 5.9, é apresentada a distribuição das questões, levando -se em conta, simultaneamente, a classificação dos dois índices.
Os dados da Tabela 5.9 atestam que sete questões que tiveram índice de discriminação “ Muito bom ” figuraram em três níveis de dificuldade – “ Difícil, Fácil e Muito Fácil ” : uma classificada na categoria “ Muito Fácil ” ( questão 6 ) ; cinco classificadas na categoria “ Fácil ” ( questões 1, 3, 4, 5 e 8 ) ; e uma, na categoria “ Difícil ” ( questão 9 ).
Em particular, a questão 4 foi a que apresentou o maior poder discriminatório, com índice 0,47, considerado “ Fácil ” em termos de facilidade, com uma proporção de 0,71 acertos, não tendo essa questão alcançado o máximo de acertos.
No Gráfico 5.4, para exemplificar, analisa -se o comportamento da Questão Objetiva 5 de Formação Geral.
Trata -se de uma questão considerada com nível “ Fácil ” de dificuldade, em relação à facilidade, e que obteve o índice de discriminação 0,46 ( muito bom ).
No Gráfico 5.4, cada uma das cinco curvas representa o percentual de respostas em determinada alternativa da questão, em função do número de acertos dos estudantes nessa parte da prova ( Formação Geral/múltipla escolha ), antes de possíveis eliminações pelo critério ponto-bisserial.
A curva em lilás corresponde à alternativa E, a alternativa correta para essa questão.
Observa -se que, entre os estudantes com menor número de acertos, nessa parte do exame, a situação mais frequente foi a escolha da alternativa C ( em verde ), incorreta.
Por exemplo, entre os estudantes que acertaram até duas respostas, 22,7 \% escolheram a alternativa A ( em vermelho ) ; 16,0 \% escolheram a alternativa B ( em mostarda ) ; 31,4 \% escolheram a alternativa C ( em verde ) ; 17,1 \% escolheram a alternativa D ( em azul, gabarito ) ; e 10,7 \% escolheram a alternativa E ( em lilás ).
Entre os que acertaram duas respostas entre as questões de múltipla escolha de Formação Geral, 1,2 \% deixou a questão em branco e 0,9 \% marcou mais de uma alternativa, invalidando a questão.
À medida que o número de acertos aumenta, indicando desempenho melhor nessa parte da prova, aumenta, concomitantemente, a proporção de estudantes que selecionaram a alternativa correta E, atingindo 100 \% para os estudantes com oito acertos.
Essa análise permite verificar como a questão discriminou os grupos de desempenho, justificando -se o alto índice ( 0,46 ) obtido na questão.
Os gráficos relativos às demais questões de Formação Geral constam do Anexo IV. 5.2.2 COMPONENTE DE CONHECIMENTO ESPECÍFICO Na Tabela 5.10, são apresentadas as estatísticas básicas em relação às questões objetivas do componente de Conhecimento Específico da prova de Medicina por grande região.
Em relação ao componente de Conhecimento Específico da prova de Medicina, como será discutido ainda nesta seção, duas questões não foram usadas no cômputo das notas por terem sido descartadas pelo critério do ponto-bisserial.
Assim, as notas foram calculadas com base em 27 das 29 questões objetivas válidas do componente de Conhecimento Específico.
Como se verifica a partir dos dados da Tabela 5.10, a média do Brasil para o componente de Conhecimento Específico da prova foi 65,2.
A menor média foi observada na região Norte ( 62,2 ), e a maior, na região Sul ( 67,7 ).
O desvio-padrão de todo o Brasil foi 14,0, sendo o maior desvio-padrão encontrado na região Sudeste ( 14,1 ), e o menor, na região Norte ( 13,3 ).
A mediana de todo o Brasil foi 66,7, o mesmo valor da mediana encontrada nas regiões Centro-Oeste, Sudeste e Sul, e superior às medianas das regiões Nordeste e Norte ( 63,0 ).
A nota máxima da prova foi 100,0, obtida, nas questões objetivas do componente de Conhecimento Específico, por, pelo menos, um estudante das regiões Centro-Oeste, Nordeste, Sudeste e Sul, na região Norte, a nota máxima foi 96,3.
A nota mínima de maior valor foi 3,7, obtida por, pelo menos, um estudante na região Nordeste.
Na Tabela 5.11, são apresentados os índices de facilidade e de discriminação ( ponto-bisserial ) das questões objetivas do componente de Conhecimento Específico, para os estudantes de Medicina.
A partir do índice de facilidade obtido, apresentado na Tabela 5.11, pode -se concluir que a maioria das questões objetivas da prova foram classificadas no índice “ Fácil ” ( 16 ).
Cinco questões foram classificadas no índice"Difícil ” ; seis, no índice “ Médio ” ; duas, no índice “ Muito fácil ” ; e nenhuma questão foi \textbf{alocada} no índice “ Muito difícil ”.
Quanto aos índices de discriminação das questões objetivas do componente de Conhecimento Específico da prova, obtém -se como resultado a seguinte classificação : quatro questões foram \textbf{alocadas} no índice “ Muito bom ” e dez, no “ Bom ”.
Dessa forma, 14 das 29 questões foram classificadas com os índices de discriminação “ Bom ” e “ Muito bom ”.
Entre as demais, 13 delas foram \textbf{alocadas} no índice “ Médio ” e duas, no “ Fraco ”.
Constata -se, assim, que a prova — no que se refere ao componente de Conhecimento Específico — possuía média capacidade de discriminar entre estudantes que dominam ou não o conteúdo.
O índice de facilidade variou de 0,16 a 0,96, e o de discriminação, de 0,03 a 0,44.
Na Tabela 5.12, é apresentada a distribuição das questões considerando -se, simultaneamente, a classificação dos dois índices.
Os dados das tabelas 5.11 e 5.12 atestam que dez questões que tiveram índice de discriminação “ Bom ” figuraram em duas níveis de dificuldade — “ Fácil"e"Médio ” : oito classificadas na categoria “ Fácil ” ; duas, na categoria “ Médio ” ; nenhuma, na categoria “ Difícil ” ; nenhuma, na categoria “ Muito fácil ” e nenhum, na categoria “ Muito difícil ”, Apenas questões tiveram o índice de discriminação “ Fraco ” e configuraram -se com índices de facilidade “ Difícil ”.
Entre as questões que alcançaram os maiores índices de discriminação, “ Muito bom ”, uma delas, a questão 36, teve o maior índice de discriminação ( 0,44 ).
Tal questão foi classificada como “ Médio ”, com índice de facilidade igual a 0,55.
A questão 20 foi a mais “ Difícil ” entre as 29 questões específicas válidas, com baixo índice de facilidade, apenas 0,16 de proporção de acertos.
Essa questão apresentou poder discriminatório médio, 0,20, o que comprova ter sido a mais “ Difícil ” para os estudantes.
Destacam -se, também, com baixo índice de facilidade, as questões 21 e 20, cujos índices de Destaca -se que as questões 21 e 29 foram consideradas inadequadas pelo critério ponto-bisserial.
Por isso, foram eliminadas do cômputo da nota final.
A título de exemplo das análises do comportamento das questões objetivas, no Gráfico 5.5, analisa -se a questão 36 do componente de Conhecimento Específico.
No Gráfico 5.5, cada uma das cinco curvas representam o percentual de respostas em determinada alternativa da questão, em função do número de acertos dos estudantes nessa parte da prova ( Componente Específico ), antes de possíveis eliminações pelo critério ponto-bisserial.
A alternativa correta B, representada no gráfico pela curva em verde, foi escolhida em maiores proporções pelos estudantes com desempenho melhor nessa parte da prova.
Já as alternativas incorretas, também denominadas distratores, foram selecionadas, principalmente, por aqueles com notas mais baixas.
A Questão 36 foi considerada, pelas respostas dos estudantes avaliados na prova, como uma questão com índice de facilidade “ Médio, 0,55 ”, ou seja, 55,0 \% dos estudantes assinalaram, acertadamente, a opção B correspondente ao gabarito.
Seu índice de discriminação foi igual a 0,44, classificado como “ Muito bom ”, o maior valor de discriminação.
Observa -se que a soma não é 100 \%, em razão das questões não respondidas ou com mais de uma opção marcada.
Cerca de 38,5 \% daqueles com nota zero marcaram a opção A ( linha vermelho ) e D ( linha vermelho ).
A proporção de estudantes que selecionaram a resposta correta B aumenta gradativamente, chegando a atingir 100 \% para 22 acertos, enquanto a proporção dos que escolheram alternativas incorretas decai, a partir de sete acertos, em função do número de acertos nessa parte da prova.
Os gráficos relativos às demais questões do componente de Conhecimento Específico constam no Anexo IV. 5.3 ANÁLISE DAS QUESTÕES DISCURSIVAS Esta seção apresenta estatísticas selecionadas e histogramas das questões discursivas de Formação Geral ( 5.3.1 ) e de Conhecimento Específico ( 5.3.2 ).
Expõe e compara, também, as médias de subpopulações, caracterizadas por grande região.
Cumpre notar que nem todas as questões passam pelo crivo dos corretores.
Após a digitalização, as respostas passam por um pré-processamento de identificação automático.
Algumas dessas questões são definidas como um “ branco automático ” ( que não são enviadas para os professores para correção, a não ser que, na segunda fase, a de inspeção visual, seja constatada uma resposta ).
O critério usado nesse pré-processamento consiste em verificar, para cada questão, a partir da quantidade de pixels, o que, provavelmente, é uma resposta em branco.
A partir da separação, os documentos que foram considerados em branco são enviados para uma verificação visual, feita por uma equipe de colaboradores devidamente capacitados.
Desse modo, qualquer erro, presente no pré-processamento, é encontrado, e as questões preenchidas são enviadas para a correção dos avaliadores.
Eventualmente, algumas respostas em branco não são detectadas nesse pré-processamento automático.
A presença de pequenos riscos, ou manchas, pode fazer o programa não identificar uma questão como um possível branco.
Essas são enviadas para a correção, mas são classificadas como “ branco ” pelos avaliadores, no momento da correção. 5.3.1 COMPONENTE DE FORMAÇÃO GERAL As análises dos resultados de desempenho dos estudantes de Medicina, na questão discursiva relativa à Formação Geral, encontram -se na Tabela 5.13 e no Gráfico 5.6, cujos dados são analisados de forma sequenciada.
No caso da Tabela 5.13, analisa -se a média das notas da questão discursiva de Formação Geral, que leva em conta as correções realizadas por duas bancas : uma que analisa e pontua aspectos relativos ao conteúdo das respostas, seguindo o padrão de respostas divulgado pelo Inep ( ver Anexo VIII ) ; e outra que avalia o desempenho linguístico expresso pela redação das respostas, seguindo um padrão de respostas específico da área de Língua Portuguesa, detalhado na seção 5.3.1.5 deste relatório.
No cálculo da nota da questão discursiva, consideram -se 80 \% relativos à nota da correção de conteúdo e 20 \% referentes à nota de desempenho linguístico.
Na Tabela 5.13, observa -se que os estudantes de todo o Brasil obtiveram, em Formação Geral, média de 58,3 nas questões discursivas.
Pode -se notar, também, que o desvio-padrão nesse conjunto de questões foi de 22,9.
A maior média foi obtida na região Sul ( 59,4 ), e a menor, na região Norte ( 57,1 ).
A mediana de todo o Brasil, neste componente, foi 59,5 e, na região Sul ( 60,1 ), foi encontrada a maior mediana.
A nota máxima ( 100,0 ) foi obtida por, pelo menos, um estudante da região Centro-Oeste, Nordeste, Norte, Sudeste e Sul.
A nota mínima ( 0,0 ) foi a mesma em todas as regiões do Brasil.
No Gráfico 5.6, está representada a distribuição das notas na questão discursiva do componente de Formação Geral.
Como mostra o Gráfico 5.6, a moda da distribuição ocorre no intervalo ( 50 ; 60 ], com frequência de 16,6 \% dos participantes, sendo que, no intervalo [ 0 ; 10 ], inclui -se, além da nota zero, a frequência de estudantes que deixaram esse tipo de questão em branco.
Na sequência, os resultados verificados para a questão discursiva de Formação Geral serão apresentados, estabelecendo -se relações com os conteúdos abordados em cada uma delas.
Os comentários das bancas de docentes avaliadores, a respeito do observado na correção das respostas dos estudantes, suas impressões e conclusões serão apresentadas junto à análise de cada questão.
Cumpre esclarecer que, tendo em vista que a questão discursiva de Formação Geral foi padronizada, ou seja, foi a mesma em todas as provas, os comentários da banca são os mesmos para todas as carreiras acadêmicas, sendo direcionados a todos os estudantes que participaram do Enade 2023.
A seguir, serão analisados o desempenho linguístico dos estudantes da área de Medicina na questão discursiva de Formação Geral do Enade 2023, apresentando -se os resultados obtidos e os comentários da banca de correção sobre a questão. 5.3.1.1 ANÁLISE DE CONTEÚDO DA QUESTÃO DISCURSIVA DO COMPONENTE DE FORMAÇÃO GERAL Os dados de Medicina, obtidos a partir das respostas à questão discursiva do componente de Formação Geral, encontram -se na Tabela 5.14 e no Gráfico 5.7, exposto após a apresentação dos dados da referida tabela.
Verifica -se, a partir dos dados da Tabela 5.14, que, no conteúdo da questão discursiva do componente de Formação Geral, os estudantes de todo o Brasil obtiveram média de 54,5.
A maior média foi obtida na região Sul ( 55,6 ), e a menor, na região Nordeste ( 53,3 ).
Quanto à variabilidade das notas, o desvio-padrão de todo o Brasil foi 27,3.
O menor desvio-padrão foi obtido na região Sudeste ( 27,0 ), e o maior desvio-padrão foi obtido pela região Nordeste ( 27,7 ).
A mediana do Brasil, como um todo, foi 55,5, sendo igual, superior e superior à da região Centro-Oeste ( 55,5 ), Nordeste ( 54,5 ) e Norte ( 54,5 ) e inferior à mediana da região Sudeste e Sul ( 57,0 ).
A nota máxima da questão discursiva foi 100,0, obtida por, pelo menos, um estudante em todas as regiões.
Já a nota mínima para todas as regiões foi zero.
No Gráfico 5.7, apresenta -se a distribuição das notas de conteúdo da questão discursiva do componente de Formação Geral.
O percentual de participantes que deixaram a prova em branco ou zeraram a questão discursiva na avaliação de conteúdo foi 5,8 \%.
O Gráfico 5.7 mostra que a moda dessa distribuição ocorre no intervalo ( 40 ; 50 ], com uma frequência de 14,1 \% dos participantes. 5.3.1.2 COMENTÁRIOS SOBRE A CORREÇÃO DE CONTEÚDO DAS RESPOSTAS À QUESTÃO DISCURSIVA A questão discursiva do componente de Formação Geral foi composta por um texto motivador e por dois itens ( a e b ).
O texto motivador apresentou um comentário sobre uma publicação do Instituto Brasileiro de Geografia e Estatística ( IBGE ) que tratava dos indicadores sobre a realidade social brasileira, apontando o crescimento do nível de desigualdade socioeconômica em 2021 e concluindo que a pobreza teve aumento recorde nesse período.
A questão foi dividida em dois itens.
O primeiro item solicitava que o respondente estabelecesse uma relação entre o perfil da população mais afetada pelas desigualdades sociais e as situações de risco socioambiental.
Já o segundo item pedia a proposição de duas possíveis ações a serem desenvolvidas em bairros periféricos a fim de minimizar os impactos socioambientais.
A questão foi pertinente, apresentou linguagem acessível e abordou conteúdo basilar para Formação Geral ; no entanto, apresentou um nível de dificuldade alto, dadas as competências intrincadas que demandou.
Além disso, o tema foi atual e de grande relevância para a sociedade brasileira, visto que abordou assunto de impacto social e de provável domínio dos concluintes de curso superior.
O texto motivador foi adequado e coerente com o que se pediu na questão, tratou do aumento da desigualdade social e da descrição do perfil da população com base em critérios socioeconômicos.
Com relação às habilidades cognitivas, observou -se que as seguintes foram necessárias para responder à questão proposta : capacidade de expressar ideias de maneira clara, coesa e argumentativa por meio da produção de textos ; capacidade de inter-relacionar objetos de conhecimento nas diferentes áreas ; e capacidade de criar estratégias de ação, selecionar métodos e formular propostas.
Entre os saberes necessários para o desenvolvimento das respostas, os respondentes deveriam conhecer o meio social e/ou geográfico da população brasileira de perfil socioeconômico pobre ou extremamente pobre e deveriam saber relacionar esse meio com situações de risco socioambiental que afligem principalmente esse perfil populacional.
Portanto, deveriam, também, conhecer uma variedade de riscos socioambientais a que esse perfil populacional está sujeito, por conta do meio social e/ou geográfico em que vive.
A questão, em seu item b ), por sua vez, demandou o domínio de um rol de propostas minimamente \textbf{factíveis} para reduzir tais situações de risco, envolvendo distintos atores sociais.
Em relação ao padrão de resposta, esperava -se para cada um dos itens avaliados que : a ) O estudante explicasse a relação entre desigualdade social e riscos socioambientais percorrendo o seguinte trajeto teórico-argumentativo : • descrever o(s) risco(s) socioambiental(is) que afeta(m) a população com perfil pobre ou extremamente pobre, tais como inundações, \textbf{deslizamentos} de terra, contaminação ambiental, maior vulnerabilidade em relação a doenças, discriminação social e criminalidade, entre outros ; • descrever o meio ( geográfico ou social ), como morros, áreas próximas a rios, mangues, espaços marcados pela violência urbana, entre outros, nas cidades brasileiras, como fator de criação ou potencialização das condições de vulnerabilidade e risco aos pobres ou extremamente pobres ; • estabelecer nexo causal entre a pobreza ou extrema pobreza, o meio ( geográfico ou social ) urbano e a situação de risco socioambiental, relacionando esses aspectos ao perfil socioeconômico da população.
Portanto, esperava -se que o estudante apresentasse consistência argumentativa quanto ao que foi solicitado, notadamente quanto aos riscos socioambientais, à descrição do meio e ao estabelecimento do nexo causal entre pobreza, meio e risco socioambiental. b ) O estudante apresentasse propostas pertinentes, \textbf{factíveis} e bem desenvolvidas que envolvessem ação governamental ( federal, estadual/distrital ou municipal ) e participação da comunidade a fim de minimizar riscos socioambientais, como, por exemplo : • implantar obras de infraestrutura urbana que envolvessem contenção de morros e encostas ; saneamento básico : serviços regulares de abastecimento de água, coleta e tratamento de esgotos, coleta e manejo de resíduos sólidos, manejo de águas \textbf{pluviais}, entre outras ; • desenvolver políticas sociais, como segurança pública, provisão de moradia adequada, com custo acessível, e regularizada do ponto de vista fundiário ; educação de qualidade ; segurança pública ; atendimento psicossocial, com especial atenção para grupos em situação de vulnerabilidade ; entre outras ;  garantir o acesso dessas comunidades a equipamentos sociais e culturais, bem como a espaços públicos inclusivos e a áreas verdes ;  elaborar planos estratégicos participativos urbanos e ambientais a serem desenvolvidos nas comunidades em situação de vulnerabilidade.
No tocante ao subitem b ), esperava -se que o respondente apresentasse coerência propositiva, englobando os diferentes atores, quais sejam : o governo, em suas diversas instâncias, e a comunidade, também em suas diversas organizações, almejando conter os riscos socioambientais. 5.3.1.5 CONSIDERAÇÕES FINAIS SOBRE A CORREÇÃO DE CONTEÚDO DA QUESTÃO DISCURSIVA DO COMPONENTE DE FORMAÇÃO GERAL De uma maneira geral, os respondentes apresentaram um nível aquém do esperado em termos de formação geral, destacando -se, neste âmbito, os aspectos a seguir.
Nas respostas relacionadas ao item a ), percebeu -se que, quanto aos riscos, a maioria dos respondentes não se limitou a descrevê -los, ou mesmo a estabelecer um rol deles.
No caso daqueles que descreveram as diversas situações de risco, em sua maioria, as respostas se enquadraram naquilo que estava previsto no padrão de resposta, apresentando os riscos derivados da insegurança relativa ao meio geográfico em que habitam, tais como alagamentos, \textbf{deslizamentos} e doenças por contato com um meio \textbf{insalubre}.
Quanto à descrição do meio, notou -se que houve uma maior aderência por parte dos respondentes, pois a maior parte das respostas, excluindo -se as que foram deixadas em branco, descreveu o meio de alguma forma.
Devido à grande variedade de possíveis descrições do meio geográfico e/ou social em que vive a população brasileira de perfil pobre ou extremamente pobre, o padrão de respostas deixava aberta a possibilidade de considerar respostas variadas em relação às citadas como exemplo em sua redação.
E, de fato, entre as respostas encontradas, verificaram -se muitas que foram consideradas certas, por estarem dentro do perfil previsto como correto, mas que eram distintas dos exemplos específicos apresentados no padrão de respostas.
O mesmo se deu quanto ao item b ) : houve uma grande variedade de proposições para além dos campos temáticos específicos previstos no padrão de respostas, mas que foram consideradas corretas porque atendiam à demanda do item.
Para esse item, o padrão de respostas demandava que as proposições fossem pertinentes ao tema e minimamente desenvolvidas.
No entanto, a maior parte dos respondentes não apresentou proposições desenvolvidas e muitos não fizeram proposições pertinentes ao tema.
O item também demandava que as proposições abrangessem a ação governamental e a participação da comunidade, mas a maior parte dos respondentes envolveu em suas proposições apenas os governos em suas mais diversas instâncias, não abrangendo o papel participativo da comunidade.
Um pequeno grupo de respostas se limitou a apresentar só uma proposição.
Dessa forma, o desempenho dos estudantes foi, em geral, muito baixo.
O número de questões deixadas em branco foi considerável, bem como o quantitativo de respostas muito curtas, ocupando menos de cinco das 15 linhas oferecidas para a resposta.
Muitos responderam a apenas um dos dois itens.
A grande maioria dos que buscaram responder ao que a questão demandava, completamente ou parcialmente, não souberam se expressar de modo correto, objetivo e coerente.
Não poucos respondentes repetiram longos trechos do enunciado em suas respostas, demonstrando um baixo nível de formação geral e, também, baixa disposição para a realização da prova.
Muitos respondentes foram desatentos e escreveram a resposta à questão discursiva de componente de Conhecimento Específico no campo dedicado à questão discursiva de Formação Geral, tendo sido incluídos dentre os respondentes que fugiram ao tema.
Uma quantidade proporcionalmente pequena, mas grande numericamente, foi a de estudantes que apresentaram respostas que desrespeitavam os direitos humanos, em sua maioria propondo o controle reprodutivo sobre a população pobre ou extremamente pobre como solução governamental para as situações de risco socioambiental comumente vivenciadas pelas populações brasileiras pobres ou extremamente pobres.
Alguns estudantes também apontaram uma relação necessária entre pobreza e \textbf{ignorância} ou mesmo entre pobreza e violência.
Entre os estudantes que responderam com maior \textbf{dedicação}, buscando dar conta de toda a questão, houve um número significativo que aparentou ter grandes dificuldade de expressão textual, dando a entender que seu baixo domínio da linguagem verbal comprometia a capacidade de produzir respostas adequadas.
Muitas respostas não apresentaram o mínimo de lógica ou não buscaram estabelecer relações entre os conteúdos mobilizados.
Portanto, em sua maioria, os estudantes não apresentaram a capacidade de expressar ideias de maneira clara, coesa e argumentativa por meio da produção de textos, uma habilidade fundamental ao concluinte do Ensino Superior.
Um outro perfil de destaque foi o de estudantes que se dedicaram a pontuar elementos válidos para responder à questão sem utilizar os conectivos, as relações causais.
Também merecem destaques os respondentes que falaram sobre aspectos históricos da pobreza e a descrição dos tipos sociais mais envolvidos em situação de pobreza ou extrema pobreza, fugindo de buscar estabelecer uma relação causal com as situações de risco socioambiental, portanto, fugindo do que a questão solicitava.
Essas situações demonstraram um perfil majoritário com baixa capacidade de produzir para o item de Formação Geral uma resposta adequada esperada para estudantes concluintes do Ensino Superior.
Os estudantes com perfis de respostas minimamente adequadas ou boas corresponderam ao segundo perfil majoritário.
Em sua maior parte, esse perfil de estudantes respondeu bem ao subitem a ), mas não respondeu bem ou subitem b ), ou vice-versa.
O perfil minoritário foi o de estudantes de nível muito bom a ótimo.
Esses com competência para descrever objetivamente os pontos essenciais da questão, estabelecendo o nexo causal entre eles.
Quanto à relação com o conteúdo, a maior parte dos respondentes soube pensar nos elementos e nas imagens do problema levantado na questão, mas não soube fazer as ligações que estimulam a relação entre esses elementos.
Portanto, demonstraram uma baixa capacidade de inter-relacionar objetos de conhecimento nas diferentes áreas, outra habilidade esperada de um concluinte do Ensino Superior.
Uma grande parte dos estudantes vê o saneamento básico precário como a maior fonte de riscos socioambientais para a população pobre ou extremamente pobre.
As descrições do meio físico, como as situações de moradias em encostas de morros e \textbf{beiras} de rio, também constituíram entre as mais comuns.
Em menor número, mas ainda com relativa frequência, apareceram os casos de estudantes que apontaram o meio como lugar de domínio territorial de criminosos ou os casos que indicaram o meio social como opressivo aos negros.
Quanto às proposições, os estudantes, em sua maioria, não souberam criar estratégias de ação, selecionar métodos e formular propostas.
Uma quantidade considerável de estudantes apresentou proposições muito genéricas como “ construir escolas ”.
A maioria não apontou conexão entre as propostas e a minimização dos riscos socioambientais, como se a ação proposta já evidenciasse o que ela está solucionando, assim também indicando uma grande concentração de estudantes que não aparenta dominar habilidades importantes para concluintes do Ensino Superior. 5.3.1.6 ANÁLISE DE LÍNGUA PORTUGUESA DA QUESTÃO DISCURSIVA DO COMPONENTE DE FORMAÇÃO GERAL Os dados de Medicina obtidos a partir das respostas à questão discursiva do componente de Formação Geral, no que tange ao desempenho linguístico, são apresentados na Tabela 5.15.
Os dados da Tabela 5.15 mostram que, no quesito desempenho linguístico, os estudantes, de todo o Brasil, obtiveram média de 73,6.
A maior média em relação ao desempenho linguístico foi obtida na região Sul ( 75,0 ), e a menor, na região Norte ( 71,8 ).
Quanto à variabilidade das notas, o desvio-padrão de todo o Brasil foi 16,2.
O menor desvio-padrão foi obtido na região Sul ( 15,3 ), e o maior desvio-padrão foi obtido na região Norte ( 17,8 ).
A mediana das notas de Língua Portuguesa foi 72,5 para o Brasil, como um todo, a mesma obtida em todas as regiões, exceto no Sul ( 77,5 ).
A nota máxima para todo o Brasil foi 100,0, com, pelo menos, um estudante com essa nota em todas as regiões.
A nota mínima foi 0,0 em todas as regiões do país.
No Gráfico 5.8, mostra -se a distribuição das notas de desempenho linguístico do componente de Formação Geral.
O Gráfico 5.8 demonstra que 1,2 \% dos participantes correspondem ao intervalo [ 0 ; 10 ].
A moda dessa distribuição ocorre no intervalo ( 70 ; 80 ], com uma frequência de 31,2 \% dos participantes. 5.3.1.7 COMENTÁRIOS SOBRE A CORREÇÃO DE LÍNGUA PORTUGUESA DAS RESPOSTAS À QUESTÃO DISCURSIVA Os comentários feitos neste item a respeito da correção de Língua Portuguesa das respostas emitidas para a questão discursiva do componente de Formação Geral estão agrupados em três pontos : avaliação da questão, padrão de resposta utilizado na avaliação e análise das respostas dadas pelos estudantes.
A linguagem da questão e do texto motivador foi considerada condizente com o domínio linguístico que se espera de estudantes concluintes de cursos em nível de graduação, sendo -lhes, portanto, acessível.
Do ponto de vista linguístico e textual, o nível da questão foi considerado mediano.
Por meio de um comando claro e objetivo, a questão solicitava que o estudante explicasse determinada relação, exposta no texto motivador, e apresentasse duas propostas para \textbf{atenuar} o problema discutido.
Para compreender o que deveria ser feito, o estudante necessitava ter domínio da leitura e interpretação textual, de modo a entender o enunciado, estabelecer a relação proposta na questão e apontar possíveis ações governamentais e da comunidade para resolver os problemas elencados.
O fato de a proposta textual ser apresentada em dois itens ( a e b ) separadamente levou muitos estudantes a produzirem o texto em dois parágrafos bem definidos e, na maioria dos casos, sem conexão textual entre ambos, ou seja, respondendo à questão em dois ou mais parágrafos que indicavam a letra do item ao qual a resposta se dirigia.
Foram relativamente poucas as provas em que os estudantes respondiam à questão em um texto que articulasse textualmente os dois itens.
Essa formatação dificultou a avaliação do quesito Aspectos Textuais no que se refere à organização textual.
Como a questão demandava certa capacidade de análise, a resposta exigia, em termos textuais, que o estudante fosse além do texto motivador, o que, em parte, passou a identificar a qualidade textual da resposta, ainda que tomasse como base o próprio enunciado da questão e empregasse dados do texto motivador.
Textos com qualidade mediana, em geral, consistiram na retextualização da questão e/ou de enunciados do texto motivador, com alguma autonomia linguística para elaborar argumentação superficial.
A simples cópia — às vezes, com erros ortográficos ou morfossintáticos — da questão ou de enunciados do texto motivador, praticamente sem autonomia textual ou argumentativa, foi um recurso característico de respostas de baixa qualidade em termos linguísticos, porém não penalizado na avaliação do desempenho linguístico.
Em relação ao padrão de resposta utilizado na avaliação, foram considerados os seguintes aspectos : ( 1 ) aspectos ortográficos ; ( 2 ) aspectos textuais ; ( 3 ) aspectos morfossintáticos e de propriedade vocabular adequada à modalidade escrita formal da Língua Portuguesa, exigida pela situação comunicativa.
A seguir, são apresentados cada um dos aspectos avaliados : 1 ) Aspectos ortográficos – referem -se ao domínio das convenções ortográficas da norma-padrão da Língua Portuguesa, como : grafia de vogais e consoantes, uso de maiúsculas e minúsculas, emprego do hífen e da acentuação gráfica, segmentação de palavras. 2 ) Aspectos textuais – referem -se ao domínio dos procedimentos de estruturação textual : organização interna dos períodos, emprego de conectores para a articulação lógica entre os períodos e entre os parágrafos, emprego de marcas de referenciação lexical e pronominal ; e utilização dos sinais de pontuação que contribuem para a organização lógica da frase. 3 ) Aspectos morfossintáticos e de propriedade vocabular – referem -se ao domínio das regras de caráter morfossintático estabelecidas como modelares do ponto de vista da modalidade escrita formal da norma-padrão da Língua Portuguesa : concordância nominal e verbal, regência nominal e verbal, colocação pronominal, flexão nominal e verbal, correlação entre tempos e modos verbais, ausência de marcas de oralidade.
Outra competência combinada a esta na análise das respostas foi a seleção vocabular adequada à modalidade escrita formal da Língua Portuguesa, tendo em vista a interseção entre as duas do ponto de vista das exigências do registro formal da modalidade escrita da norma-padrão.
No que concerne à análise das respostas dadas pelos estudantes, especificamente, em relação ao desempenho linguístico, foram observados muitos deslizes gramaticais no que diz respeito à acentuação e grafia correta das palavras, embora, de maneira geral, os estudantes tenham demonstrado domínio da ortografia vigente.
Quanto aos aspectos textuais, verificou -se predominância de textos organizados sem a presença das três partes esperadas para um texto dissertativo-argumentativo : introdução ; desenvolvimento e conclusão.
A argumentação foi predominantemente rasa e pouco elaborada, havendo frequentes desvios de pontuação, nomeadamente quanto ao emprego da vírgula.
Por fim, quanto aos aspectos morfossintáticos, observou -se um bom domínio das regras de regência e concordância na maior parte dos textos.
A adequação lexical também foi pertinente a uma redação formal na maior parte das vezes.
Em relação ao item a ), destacaram -se como respostas a relação entre comunidades negras periféricas, além de ribeirinhos e indivíduos em vulnerabilidade socioeconômica como sendo o perfil da população brasileira mais propenso a riscos socioambientais.
No que diz respeito ao item b ), foram encontradas com maior frequência propostas relacionadas a atuações governamentais, familiares e sociais diversas, como a necessidade de maiores investimentos e ofertas de trabalho na periferia.
Houve a ocorrência de respostas fracas e incompletas.
O item a ) foi respondido, na maior parte dos textos, de modo discursivo, enquanto o item b ) figurou, por vezes, como uma lista de ações que poderiam resolver o problema.
Em alguns casos, não tão raros, foram produzidos textos nos quais os estudantes se recusaram explicitamente a produzir uma resposta adequada, alegando que a questão era pertinente a uma área de formação diferente da sua.
As respostas encontradas e mesmo a ausência de respostas nos permitem tecer algumas reflexões.
Há constante falta de atenção na leitura do enunciado da questão, o que levou alguns estudantes a responderem à questão específica em local diverso do esperado, prejudicando a resposta.
Por vezes, alguns candidatos produziram as respostas para as duas questões ( Formação Geral e Conhecimento Específico ) no mesmo espaço, prejudicando o desempenho textual e argumentativo.
Seja por desatenção, seja por descuido, a dificuldade de interpretação é sintomática e deve ser considerada.
Desvios ortográficos, morfossintáticos e de pontuação frequentes apontam para lacunas no que concerne à alfabetização e ao desenvolvimento da escrita ainda nos anos do ensino fundamental.
Contudo, o número baixo de textos que tenham obtido notas muito baixas, como 1,0 e 2,0, revela que, entre os estudantes avaliados, não havia pessoas com níveis extremamente baixos de alfabetização e letramento.
Problemas quanto à organização do texto dissertativo-argumentativo, com ausência de introdução, desenvolvimento e conclusão, mesmo que minimamente, apontam para a necessidade de um trabalho mais profundo com a escrita em cursos superiores.
A argumentação rasa ou mesmo o desinteresse por temas sociais revelam o fato de que alguns dos cursos estudados não trabalham o senso crítico do estudante, não o conduzindo a uma vivência plenamente crítica e democrática em sociedade.
A dificuldade em responder ao item b ) da questão demonstrou uma fragilidade teórica na formação, principalmente no que concerne à atuação do profissional em relação a questões de cunho socioeconômico e ao atendimento ao público mais fragilizado socioeconomicamente.
A baixa frequência de desvios morfossintáticos pode ser justificada pela própria competência linguística do estudante, demonstrando que um trabalho mínimo com questões gramaticais nos cursos superiores poderia sanar grande parte do problema.
Relativamente poucos estudantes demonstraram incapacidade quanto à escolha de palavras adequadas para a produção de um texto formal, apresentando competência genérica, na maior parte das vezes, satisfatória.
Textos que obtiveram nota máxima foram raros, o que aponta para a necessidade de reforçar a importância do trabalho com a leitura e escrita em âmbito formal com os estudantes dos cursos avaliados. 5.3.1.8 CONSIDERAÇÕES FINAIS SOBRE A CORREÇÃO DE LÍNGUA PORTUGUESA DAS RESPOSTAS À QUESTÃO DISCURSIVA DO COMPONENTE DE FORMAÇÃO GERAL A análise global em relação ao desempenho linguístico dos cursos avaliados no Enade 2023 revelou um desempenho mediano dos estudantes avaliados.
Quanto aos aspectos textuais e argumentativos, verificou -se um nível baixo de conhecimento e mesmo de interesse, por parte de muitos estudantes, em refletir sobre os problemas sociais envolvendo pessoas em situação de vulnerabilidade e fenômenos de risco socioambiental.
A maioria dos estudantes não conseguiu apresentar uma análise crítica, tampouco responder conforme solicitado pela questão.
Alguns estudantes não foram atentos às informações apresentadas nos textos nem ao que os enunciados solicitavam.
Alguns demonstraram o domínio sobre o assunto ou, no mínimo, o comprometimento e o interesse para responder as questões.
Entende -se que o tempo, a quantidade de questões e a pressão podem influenciar no momento de elaborar a resposta.
Além disso, alguns estudantes de cursos específicos não se sentiram motivados a falar de um assunto julgado como “ social ” e não pertinente à sua área de atuação profissional, o que é um forte representativo de vieses ideológicos que perpassam determinadas formações.
Cabe ressaltar que a questão abordava um conteúdo atual e relevante a ser considerado por um indivíduo plenamente capaz de agir em sociedade de forma crítica, de qualquer área de formação em nível superior.
Destaca -se que estudantes com aparente domínio da área de saúde apresentaram inclinação crítica a abordar o tema de forma bastante humanizada.
Quanto aos aspectos ortográficos e gramaticais de natureza morfossintática, observam -se muitas lacunas trazidas desde a educação básica.
Entende -se que a universidade não pode sanar todos os problemas de cunho ortográfico e gramatical, mas é preciso salientar a importância de disciplinas voltadas à leitura e produção de texto serem incluídas como disciplinas obrigatórias nos cursos avaliados. 5.3.2 COMPONENTE DE CONHECIMENTO ESPECÍFICO As análises dos resultados de desempenho dos estudantes de Medicina, na questão discursiva relativa ao componente de Conhecimento Específico, encontram -se na Tabela 5.16 e no Gráfico 5.9. 5.3.2.1 ANÁLISE DA QUESTÃO DISCURSIVA DO COMPONENTE DE CONHECIMENTO ESPECÍFICO A Tabela 5.16 apresenta dados sobre a questão discursiva do componente de Conhecimento Específico.
A partir dos dados registrados na Tabela 5.16, observa -se que a média foi mais baixa se comparada à média obtida na questão discursiva do componente de Formação Geral.
Enquanto, no componente de Formação Geral, a média para estudantes de Medicina de todo o Brasil foi 58,3 ; na parte de Conhecimento Específico, a média foi 60,0.
A maior média desse componente foi obtida pelos estudantes da região Sul ( 63,3 ), e a menor, pelos da região Norte ( 57,4 ).
Quanto à variabilidade das notas, o desvio-padrão de todo o Brasil foi 20,7.
O maior desvio-padrão foi encontrado nas regiões Nordeste ( 21,2 ), e o menor, na região Sul ( 19,8 ).
A nota máxima obtida foi 100,0 e a nota mínima foi 0,0 em todas as regiões do país.
No Gráfico 5.9, apresenta -se a distribuição das notas na questão discursiva no componente de Conhecimento Específico.
À luz dos dados do Gráfico 5.9, verifica -se que a moda da distribuição das notas da questão discursiva do componente de Conhecimento Específico ocorre no intervalo ( 60 ; 70 ], com 17,8 \% do total de participantes.
Observa -se um valor considerável no intervalo [ 0 ; 10 ], cerca de 1,2 \% dos participantes, entretanto, vale ressaltar que, nesse intervalo, são contabilizadas também as questões deixadas em branco. 5.3.2.2 COMENTÁRIOS SOBRE AS RESPOSTAS À QUESTÃO DISCURSIVA A questão discursiva apresentada para área de Medicina foi composta por um caso clínico e três itens ( a, b e c ).
O caso clínico apresentado referia -se a uma paciente hipotética atendida em uma Unidade Básica de Saúde, com três diagnósticos principais : diabetes mellitus, hipertensão arterial sistêmica e dislipidemia.
Os itens da questão ( a, b e c ) exploraram o conhecimento do estudante em relação à abordagem terapêutica para os problemas principais apresentados, assim como orientações gerais e respostas a dúvidas específicas da paciente.
O caso clínico apresentado no texto abordou uma situação frequente na prática do estudante de Medicina, que certamente foi vivenciada durante sua formação.
Foi apresentada uma mulher preta, com 58 anos de idade, atendida na Unidade Básica de Saúde, que consumia bebida alcoólica e apresentava sobrepeso.
Os aspetos macroestruturais explorados foram desde considerações primordialmente técnico-científicas até aquelas relativas ao relacionamento interpessoal, explorando a empatia e a capacidade de dialogar e de se expressar utilizando uma linguagem simples e de fácil entendimento pelo paciente, o que é crucial para a formação do médico.
Nas respostas dadas pelos estudantes, foi frequente a inobservância quanto ao detalhamento de alguns itens, conforme exigido no padrão de respostas.
As respostas denotaram que os estudantes, muitas vezes, tinham o conhecimento em relação ao tema, mas não entenderam a necessidade de aprofundar a explicação.
Como exemplo, no item a ), foi frequente a indicação da atividade física, mas com pouca frequência os estudantes recomendaram a perda ponderal como medida adicional para o tratamento das comorbidades.
Frequentemente, foram detalhados outros aspectos da abordagem não farmacológica não contemplados no padrão de respostas, como, por exemplo, a melhora na qualidade do sono, o encaminhamento para nutricionista e o suporte psicológico.
O comando do item b ) exigia a apresentação completa do plano terapêutico.
Entretanto, muitos não entenderam a necessidade de indicar os nomes farmacológicos, respondendo ao item com nomes comerciais de medicamentos, talvez em virtude de seu emprego rotineiro na atividade prática dos estudantes.
Foi frequente a prescrição equivocada de alguns grupos anti-hipertensivos, como diuréticos de alça, betabloqueadores e insulina.
Alguns estudantes não entenderam que a resposta ao item c ) deveria ser elaborada em linguagem adequada para o entendimento pela paciente, de forma a esclarecer suas dúvidas.
Muitos mantiveram uma linguagem extremamente técnica, dificultando o entendimento da paciente.
Muitos estudantes também informaram sobre a alteração da função renal, mas sem relacioná -la ao diabetes e à hipertensão, perdendo, portanto, pontuação no item.
No item a ), esperava -se que o estudante elaborasse, em sua resposta, um plano terapêutico contendo dieta e outras duas medidas para o manejo não farmacológico das afecções da paciente, quais sejam diabetes mellitus, hipertensão arterial e dislipidemia.
Os temas foram abordados de forma adequada em diretrizes e/ou protocolos brasileiros para o manejo de hipertensão, diabetes e dislipidemia.
Os estudantes trouxeram respostas diversas, tendo em vista a gama enorme de alternativas nas chamadas medidas não farmacológicas na abordagem das afecções.
Frequentemente, não houve necessariamente uma hierarquia na indicação das condutas.
Nem sempre a dieta hipossódica foi lembrada, nem tampouco a cessação ou interrupção do consumo de alcoólicos.
Isso se deve, talvez, a pouca valorização da indicação da mudança de hábitos de vida em prol da medicalização nos tempos atuais.
No item b ), oram frequentemente indicados fármacos não disponíveis na rede pública ( inibidores da SGLT2, liraglutida ), talvez pelo acesso de boa parte dos estudantes ao atendimento em rede privada de saúde.
Da mesma forma, foi frequente a indicação de insulina de forma inadequada no tratamento do diabetes na condição atual da paciente apresentada no caso clínico, assim como a prescrição de fármacos de forma equivocada para a hipertensão da paciente, como por exemplo espironolactona e furosemida.
Isso se justifica, possivelmente, pelo despreparo do estudante na abordagem ao paciente com diabetes e hipertensão.
Já no item c ), muitos estudantes utilizaram termos como"comorbidades"ou"doenças de base", sem especificar diabetes e hipertensão, o que foi considerado inadequado, pois, dessa forma, a paciente poderia, por exemplo, não entender do que se tratava a comorbidade.
Nesse item, o que foi solicitado era informar a paciente, com ênfase em quais são os aspectos importantes no controle do diabetes e da hipertensão.
Tratou -se de uma questão de dificuldade mediana, que exigia conhecimentos técnico-científicos e preparo no relacionamento interpessoal, explorando a empatia e a capacidade de dialogar e de se expressar, utilizando uma linguagem simples e de fácil entendimento com a paciente, o que é crucial para a formação do médico.
Em relação a esse ponto de vista, as respostas podem ter variado de acordo com o arsenal terapêutico disponível na saúde pública e a maior ou menor exposição do estudante a atendimentos predominantemente no Sistema Único de Saúde.
Observou -se um baixo número de respostas em branco ou anuladas, assim como nota zero em baixa frequência, o que denota a adequação do problema apresentado a prevalência das situações clínicas relevantes.
As respostas encontradas e mesmo a ausência de respostas permitiram algumas reflexões que remetem a eventuais lacunas na formação médica, ou até mesmo dificuldades na interpretação dos enunciados dos itens da questão discursiva.
A maioria dos estudantes explorou de forma adequada os aspectos das medidas não farmacológicas na intervenção terapêutica para as doenças principais, evidenciando um alto índice de acertos nessa parte.
Entretanto, foram observados erros de interpretação do enunciado da questão, como a menção equivocada da cessação do etilismo como um aspecto da dieta.
Além disso, aspectos importantes, como a perda ponderal e a redução da ingestão de sal, não foram mencionados pela maioria dos estudantes.
Por outro lado, a recomendação de atividade física foi amplamente destacada nas respostas, demonstrando sua prevalência como medida não farmacológica essencial.
A falta de menção à perda ponderal e à redução da ingestão de sal por parte da maioria dos estudantes pode ser atribuída a lacunas na formação médica, indicando a necessidade de reforçar o conhecimento desses temas durante o processo de formação.
Além disso, a confusão entre a cessação do etilismo como um aspecto da dieta pode sugerir uma possível falta de compreensão sobre as nuances entre hábitos alimentares e comportamentais.
Esses erros ressaltam a importância de uma abordagem mais aprofundada e clara no ensino das medidas não farmacológicas para o tratamento de condições médicas, a fim de garantir uma compreensão precisa e abrangente por parte dos futuros profissionais de saúde.
Essa abordagem, muitas vezes, prioriza intervenções farmacológicas em detrimento da sensibilização e educação sobre mudanças de estilo de vida.
Portanto, os erros identificados sugerem a necessidade de reformas educacionais que enfatizem uma visão mais holística da Medicina, incluindo a importância das medidas não farmacológicas e a promoção de um estilo de vida saudável como pilares fundamentais no cuidado dos pacientes.
Os estudantes tiveram muitas dificuldades em responder corretamente quanto ao plano terapêutico farmacológico.
Definitivamente, essa parte da questão discursiva foi a que trouxe mais desafio.
Poucos acertaram de forma completa o tratamento da dislipidemia, errando principalmente a dose do medicamento e o horário de tomada.
Vários indicaram equivocadamente a introdução de insulina no tratamento do diabetes, para o caso concreto em análise.
Foram cometidos equívocos frequentes na grafia de princípios ativos e indicação inadequada de alguns fármacos anti-hipertensivos.
As respostas evidenciam erros substanciais em termos de apresentações, dosagens de medicamentos e intervalos posológico.
É crucial ressaltar a responsabilidade do profissional de saúde em questões tão críticas, pois erros nesses domínios podem resultar em consequências graves.
A baixa taxa de acertos na questão relacionada à prescrição de medicamentos merece reflexão, assim como o uso de nomes comerciais no lugar de princípios ativos.
Essa discussão é de extrema relevância e deve ser destacada como um assunto essencial no cenário atual da educação médica, tendo em vista a importância de formar profissionais capacitados para tomar decisões seguras e embasadas no que diz respeito à prescrição e à administração de medicamentos, visando sempre à segurança e ao bem-estar dos pacientes.
A dificuldade em responder ao item c ) da questão pode ser indício de que os estudantes precisam aprimorar o uso de linguagem compreensível e empática com os pacientes.
A maioria dos estudantes optou por termos técnicos, como classificações de doenças e estágios da doença renal, em detrimento de uma comunicação clara e acessível ao paciente.
Essa escolha inadequada na resposta pode ter sido influenciada por uma interpretação equivocada do enunciado, possivelmente devido ao cansaço gerado pela extensão da prova como um todo.
Contudo, essa falha também pode evidenciar uma lacuna na formação médica, que muitas vezes negligencia a importância do diálogo direto com os pacientes e a \textbf{dedicação} de tempo para orientações compreensíveis.
Apenas algumas respostas conseguiram estabelecer a conexão entre o diabetes e hipertensão e o impacto no comprometimento renal, destacando a relevância de abordagens preventivas para evitar a necessidade de diálise no futuro.
Erros graves foram identificados, como a indicação de que a paciente possuía função renal normal, erroneamente desconsiderando a condição real.
Alguns poucos participantes ainda cometeram o equívoco de sugerir a diálise como intervenção, o que denota a necessidade urgente de aprimorar o entendimento dos estudantes em relação a abordagem integral e adequada de casos clínicos complexos, ressaltando a importância da comunicação empática e clara com os pacientes para uma prática médica eficaz e centrada no cuidado individualizado. 5.3.2.3 CONSIDERAÇÕES FINAIS SOBRE A QUESTÃO DISCURSIVA DO COMPONENTE ESPECÍFICO Conforme exposto anteriormente, as respostas à questão discursiva demonstraram desempenho suficiente dos estudantes que se encontram já no final da formação médica.
Aspectos relevantes foram identificados em relação à prescrição de medicamentos e à comunicação com o paciente.
Para aprimorar a formação médica e abordar as lacunas identificadas nos erros apontados, é fundamental incorporar estratégias educacionais específicas que promovam uma comunicação eficaz e empática.
É evidente a necessidade de oferecer maior treinamento para a prescrição medicamentosa, em especial no que se refere à abordagem de doenças mais prevalentes, conforme a apresentada no caso clínico descrito na questão discursiva.
Algumas sugestões incluem enfatizar a comunicação com o paciente, introduzindo treinamentos e workshops que ensinem os estudantes a se comunicarem de forma clara, empática e acessível, adaptando a linguagem técnica para garantir a compreensão do paciente.
Além disso, é importante implementar simulações de casos clínicos que enfatizem a importância da abordagem integrada das condições do paciente, incentivando a prática de uma Medicina centrada no paciente.
A inclusão de aspectos comportamentais nos estudos de casos também se mostra crucial, destacando a importância dos mesmos para a promoção da saúde, como a adesão ao tratamento e o estímulo de hábitos saudáveis.
Além disso, a cada atividade prática do curso de Medicina deve ser enfatizada e praticada a abordagem da comunicação com o paciente.
A discussão sobre ética e responsabilidade deve ser ampliada na formação médica, promovendo debates sobre esses temas em relação às decisões clínicas, ressaltando a importância da precisão e empatia.
Além disso, é importante incentivar a participação em programas que abordem a importância da comunicação eficaz na prática clínica.
Essas iniciativas visam não só corrigir os erros identificados, mas também preparar os futuros médicos para uma prática clínica mais holística, empática e centrada no paciente, garantindo uma melhor qualidade nos cuidados de saúde oferecidos.
A análise global da questão discursiva da área de Medicina revelou um desempenho suficiente dos estudantes avaliados.
Em relação aos três itens da questão específica, verificou -se mais facilidade na indicação de medidas não farmacológicas para o tratamento das comorbidades indicadas no caso clínico.
Por outro lado, verificou -se um nível baixo de conhecimento em relação a correta prescrição de medicamentos, indicando nomes farmacológicos, apresentações, doses e intervalos posológicos adequados.
A maioria dos estudantes não foi capaz de apresentar uma resposta correta em relação às angústias da paciente, não conseguindo explicar de forma adequada para ela as consequências do diabetes e hipertensão para a função renal.
Alguns poucos demonstraram domínio completo sobre o assunto.
Não ficou evidente uma falta de comprometimento ou de interesse para responder à questão.
Foi muito baixo o número de questões em branco, ou que tenham sido anuladas.
O item a ) da questão teve elevado índice de acerto e quase todos os estudantes tentaram responder a todas as perguntas do enunciado.
Cabe ressaltar, por fim, que a questão específica abordou temas relevantes e prevalentes, que precisam ser muito discutidos durante a graduação, inclusive nas atividades práticas, e, conforme o preconizado nas diretrizes do curso de Medicina.
REFERÊNCIAS BÄCKSTRÖM, M.
Higher-order factors in a five-factor personality inventory and its relation to social desirability.
European Journal of Psychological Assessment, v. 23, n. 2, p. 63-70, 2007.
Disponível em : https://psycnet.apa.org/doiLanding? doi=10.1027\%2F1015-5759.23.2.63.
Acesso em : 4 mar. 2024.
BRASIL.
Lei n.º 10.861, de 14 de abril de 2004.
Institui o Sistema Nacional de Avaliação da Educação Superior ( Sinaes ) e dá outras providências.
Diário Oficial da União.
Brasília, DF, Seção 1, p. 3-4, n. 72, 15 abr. 2004.
BRASIL.
Lei n.º 13.146, de 6 de julho de 2015.
Institui a Lei Brasileira de Inclusão da Pessoa com Deficiência ( Estatuto da Pessoa com Deficiência ).
Diário Oficial da União, Brasília, DF.
Seção 1, n. 127, p. 2, 7 jul. 2015.
BRASIL.
MINISTÉRIO DA EDUCAÇÃO ( MEC ).
Portaria Normativa n.º 41, de 20 de janeiro de 2022.
Estabelece o regulamento do Exame Nacional de Desempenho dos Estudantes – Enade, edição 2022, referente ao ano III do 6.º ciclo avaliativo, e regulariza o 7.º ciclo avaliativo, previstos pelo art. 40 da Portaria MEC nº 840, de 24 de agosto de 2018.
Diário Oficial da União.
Brasília, DF, Seção 1, p. 64, n. 15, 21 jan. 2022.
CHRISTENSEN, A.
P. ; GARRIDO, L.
E. ; GOLINO, H.
Unique variable analysis : A network psychometrics method to detect local dependence.
Multivariate Behavioral Research, v. 58, n. 6, p. 1165-1182, 2023.
Disponível em : https://doi.org/10.1080/00273171.2023.2194606.
Acesso em : 5 mar. 2024.
CHRISTENSEN, A.
P. ; GOLINO, H.
On the equivalency of factor and network loadings.
Behavior Research Methods, v. 53, n. 4, p. 1563-1580, 2021.
Disponível em : https://doi.org/10.3758/s13428-020-01500-6.
Acesso em : 5 mar. 2024.
CONSELHO NACIONAL DE EDUCAÇÃO ( CNE ).
CÂMARA DE EDUCAÇÃO SUPERIOR ( CES ).
Resolução nº 5, de 14 de outubro de 2021.
Institui as Diretrizes Curriculares Nacionais do Curso de Graduação em Medicina.
Diário Oficial da União.
Brasília, DF, Seção 1, p. 47, n. 196, 18 out. 2021.
DAMÁSIO, B.
F. ; BORSA, J.
C.
Manual para desenvolvimento de instrumentos psicológicos.
São Paulo : Vetor, 2017, p. 347–374.
GOLINO, H.
F. ; EPSKAMP, S.
Exploratory graph analysis : A new approach for estimating the number of dimensions in psychological research.
PLoS ONE, v. 12, n. 6, e0174035, 2017.
Disponível em : https://doi.org/10.1371/journal.pone.0174035.
Acesso em : 6 mar. 2024.
GOLINO, H. ; SHI, D. ; CHRISTENSEN, A.
P. ; GARRIDO, L.
E. ; NIETO, M.
D. ; SADANA, R. ; THIYAGARAJAN, J.
A. ; MARTINEZ-MOLINA, A.
Investigating the performance of exploratory graph analysis and traditional techniques to identify the number of latent factors : A simulation and tutorial.
Psychological Methods, v. 25, n. 3, p. 292–320, 2020.
Disponível em : https://doi.org/10.1037/met0000255.
Acesso em : 7 mar. 2024.
HORN, J.
A rationale and test for the number of factors in factor analysis.
Psychometrika, v. 30, n. 2, p. 179-185, 1965.
Disponível em : https://doi.org/10.1007/BF02289447.
Acesso em : 5 mar. 2024.
INSTITUTO NACIONAL DE ESTUDOS E PESQUISAS EDUCACIONAIS ANÍSIO TEIXEIRA ( INEP ).
Portaria nº 109, de 1.º de abril de 2022a.
Institui as Comissões Assessoras de Área – CAA para realização de atividades referentes ao Exame Nacional de Desempenho dos Estudantes – Enade edição 2022.
Diário Oficial da União.
Brasília, DF, Seção 1, p. 3-4, n. 66-, 4 abr. 2022.
INSTITUTO NACIONAL DE ESTUDOS E PESQUISAS EDUCACIONAIS ANÍSIO TEIXEIRA ( INEP ).
Portaria nº 110, de 1.º de abril de 2022b.
Designa, com fundamento na Portaria Inep n.º 109, de 01 de abril de 2022, os membros para a constituição das Comissões Assessoras de Área ( CAA ), para realização de atividades referentes ao Exame Nacional de Desempenho dos Estudantes – edição 2022, no ano III do ciclo avaliativo previsto no art. 40 da Portaria MEC n.º 840, de 24 de agosto de 2018.
Diário Oficial da União.
Brasília, DF, Seção 2, p. 46, n. 64, 04 abr. 2022.
INSTITUTO NACIONAL DE ESTUDOS E PESQUISAS EDUCACIONAIS ANÍSIO TEIXEIRA ( INEP ).
Portaria nº 289, de 30 de junho de 2022.
Dispõe sobre diretrizes de prova e componentes específicos da área de Medicina, no âmbito do Exame Nacional de Desempenho dos Estudantes ( Enade ), edição 2022.
Diário Oficial da União.
Brasília, DF, Seção 1, p. 205, n. 123, 1.º jun. 2022.
INSTITUTO NACIONAL DE ESTUDOS E PESQUISAS EDUCACIONAIS ANÍSIO TEIXEIRA ( INEP ).
Portaria n.º 298, de 30 de junho de 2022.
Dispõe sobre diretrizes de prova e componentes específicos da área de Formação Geral, no âmbito do Exame Nacional de Desempenho dos Estudantes ( Enade ), edição 2022.
Diário Oficial da União.
Brasília, DF, Seção 1, p. 209, n. 123, 1.º jul. 2022.
INSTITUTO NACIONAL DE ESTUDOS E PESQUISAS EDUCACIONAIS ANÍSIO TEIXEIRA ( INEP ).
DIRETORIA DE AVALIAÇÃO DA EDUCAÇÃO SUPERIOR ( DAES ).
Nota Técnica Nº 07/2022/CGCQES/DAES [ 2022 ].
Apresenta a metodologia utilizada no cálculo do Conceito Enade referente ao ano de 2021, considerando os procedimentos definidos pelo Inep, conforme atribuição estabelecida pelo Decreto nº 9235, de 15 de dezembro de 2017, e pela Portaria Normativa n.º 840, de 24 de agosto de 2018, republicada em 31 de agosto de 2018, retificada em 03 de setembro de 2018.
Disponível em : https://download.inep.gov.br/educacao\_superior/enade/notas\_tecnicas/2020/nota\_tecnica\_n\_7\_2022\_CGCQES-DAES\_metodologia\_calculo\_conceito\_enade\_2021.pdf.
Acesso em 25 abr. 2023.
MANSOLF, M. ; REISE, S.
P.
Exploratory Bifactor Analysis : The Schmid-Leiman Orthogonalization and Jennrich-Bentler Analytic Rotations.
Multivariate Behavioral Research, v. 51, n. 5, p. 698-717, 2016.
Disponível em : https://doi.org/10.1080/00273171.2016.1215898.
Acesso em : 8 mar. 2024.
MEULMAN, J.
J.
Optimal scaling methods for multivariate categorical data analysis.
SPSS White Paper, 2016. 12p.
Disponível em : http://bayes.acs.unt.edu:8083/BayesContent/class/Jon/SPSS\_SC/Module9/M9\_CatReg/SWPOPT.pdf.
Acesso em : 29 abr. 2023.
PASQUALI, L.
Psicometria.
Revista da Escola de Enfermagem da USP, São Paulo, v. 43, Edição Especial, p. 992-999, dez. 2009.
PAULHUS, D.
L.
Control of social desirability in personality inventories : Principal-factor deletion.
Journal of Research in Personality, v. 15, n. 3, p. 383-388, 1981.
Disponível em : https://doi.org/10.1016/0092-6566(81)90035-0.
Acesso em : 7 mar. 2024.
PAULHUS, D.
L.
Measurement and control of response styles.
In : J.
P.
ROBINSON ; P.
R.
SHAVER ; L.
S.
WRIGHTSMAN ( Eds. ).
Measures of personality and social psychological attitudes, Academic Press, San Diego, CA, p. 17-59, 1991.
PEABODY, D.
Trait inferences : Evaluative and descriptive aspects.
Journal of Personality and Social Psychology, v. 7, n. 4, Pt. 2, p. 1-18, 1967.
Disponível em : https://doi.org/10.1037/h0025230.
Acesso em : 7 mar. 2024.
PELT, D.
H.
M. et al.
The motivation and opportunity for socially desirable responding does not alter the general factor of personality.
Assessment, v. 28, n. 5, p. 1376-1396, 2021.
PETTERSSON, E. ; TURKHEIMER, E. ; HORN, E.
E. ; MENATTI, A.
R.
The General Factor of Personality and Evaluation.
European Journal of Personality, v. 26, n. 3, p. 292-302, 2012.
Disponível em : https://doi.org/10.1002/per.839.
Acesso em : 8 mar. 2024.

% matched lemmas: alocar, atenuar, beira, dedicação, deslizamento, factível, ignorância, insalubre, pluvial
\end{textsample}
